\section{Method}
\label{sec:method}

\subsection{The production model (resolution D-05)}

XGBoost and LightGBM are trained with native handling of missing values and categories, class
imbalance addressed by weighting (no oversampling), and early stopping on the validation split.
Their raw probabilities are combined as $\nWeightXGB\,p_{\mathrm{xgb}} +
\nWeightLGB\,p_{\mathrm{lgb}}$, and \textbf{one} isotonic regression is fitted on the combined
score using the calibration split, which is chronologically later than the validation data and
never used for fitting or scoring otherwise. The calibrated ensemble score is the quantity the
decision thresholds and the calibration error are measured on; calibrating the two models
separately and averaging them does not yield a calibrated ensemble. In the evaluated run the
early-stopped models used \nGateXGBRounds{} (XGBoost) and \nGateLGBRounds{} (LightGBM) rounds,
with a positive-class weight of \nGateSPW{}.

An Isolation Forest is fitted on imputed features (it cannot take NaN) and its score is mapped to
a percentile of its training reference distribution, so that the anomaly score lies in $[0, 1]$
(resolution D-06). The specification routes transactions flagged only by the anomaly score to a
separate review queue that does not hold the payment; that routing belongs to the decision
engine (M6) and is not evaluated here.

\paragraph{Training volume.} The evaluated models were trained on \nTrainUsed{} rows holding
\nGateTrainFraud{} confirmed fraud, a subset of the \nTrainRows{} rows in the training period,
because of the development machine (Section~\ref{sec:repro}). Every figure in
Section~\ref{sec:evaluation} is at that training volume; full-volume training is
\notyet{M10}.

\subsection{Explanations}

Exact TreeSHAP is computed per model in log-odds (margin) space, where additivity is exact, and
the analyst-facing contribution is the weighted combination $\nWeightXGB\,\phi_{\mathrm{xgb}} +
\nWeightLGB\,\phi_{\mathrm{lgb}}$ in margin space, annotated with its base value and final
margin. We never claim that SHAP values sum to the calibrated probability: isotonic calibration is
monotone but not additive. The specification requires plain-language explanations to come from
deterministic, localised templates per feature and never from a language model; the templates
belong to the serving and console milestones (M5, M8).

\subsection{The diagnostic battery}

The analyses of feature groups, slices, variants and countries (Sections~\ref{sec:evaluation}
and~\ref{sec:generalisation}) use a separate, deliberately simple model: one untuned XGBoost
booster on \nTrainUsed{} training rows, calibrated with two-parameter Platt scaling. It is a
diagnostic of the benchmark, not the production model, and its figures are labelled as such
(ADR 0029). Ablations of the production ensemble are reported with the gate.

\subsection{Evaluation protocol}

\begin{itemize}
  \item \textbf{Operating points} (resolutions D-01, D-02): recall at a \nFPRBudget{}
        false-positive rate replaces the specification's precision at that rate, which is
        mathematically unreachable at this base rate; recall, F1 and the false-negative rate are
        read at the flag threshold \nFlagThreshold{}, the false-positive rate at the block
        threshold \nBlockThreshold{}.
  \item \textbf{Uncertainty}: \nResamples{} stratified bootstrap resamples for gate metrics,
        Hanley--McNeil intervals for AUCs of single features and baselines, Wilson intervals for
        recall on small variant counts, DeLong's test for AUC differences, and the mean and
        standard deviation over \nSeeds{} refits.
  \item \textbf{The floor}: every ranking metric is reported beside the best single feature and
        the best trivial rule (one threshold on one feature) on the same rows, as a margin
        (Section~\ref{sec:floor}).
  \item \textbf{Test discipline}: nothing is tuned on the test period; every read of the test set
        is logged with a counter; the calibrator is guarded by a test that fails if it sees a
        test-period row.
  \item \textbf{Scale}: every figure states the rows and fraud count it rests on, because
        single-feature AUC on this benchmark moves with dataset size by more than seed noise
        (spread \nSizeSeriesSpan{} across sizes against a seed-to-seed standard deviation of
        \nSeedSdSizeSeries{}), for a reason that remains unexplained.
\end{itemize}
